\section{Vertical ideal models of rational piecewise linear functions}

The test-configuration description of piecewise linear functions used in
Theorem~\ref{thm:wu-bj-testconfig} starts with actual ideals on the product with the
affine line.  The following construction supplies these ideals directly from
finite Fubini--Study formulas.  It concerns the ambient projective product;
its functions are restricted to the given projective subvariety.

\begin{lemma}[Integral weights and actual product ideals]
\label{lem:pl-fs-product-ideal}
\uses{def:wu-proj-functions,projective:paper-gauss-actual-product-model-log}
\lean{EnvelopeProperty.ProjBerk.FSData.exists_productIdeal_model}
\leanok
Let $D$ be rational Fubini--Study data, with
$D(v)=m^{-1}\max_j(\log|s_j|(v)+\lambda_j)$ and all $s_j$ homogeneous of
degree $d$.  There are integers $N>0$ and $C$, and a finitely generated
homogeneous ideal $J\subset k[t][x_0,\ldots,x_M]$, such that the intrinsic
model logarithm of its actual ideal sheaf on $\mathbf P^M\times\mathbf A^1$
at the Gauss lift of $v$ is
\[
 \log|\mathcal J|(\sigma(v))=mN D(v)-C.
\]
\end{lemma}
\begin{proof}
Choose $N>0$ clearing the denominators of the finitely many rational
$\lambda_j$, and write $z_j=N\lambda_j\in\mathbf Z$.  Choose an integer
$C\ge z_j$ for every $j$, and put $e_j=C-z_j\ge0$.  Take the homogeneous
ideal generated by $s_j^Nt^{e_j}$.  Each generator has $x$-degree $Nd$.
The Gauss logarithm of a generator is
$N\log|s_j|-e_j=N(\log|s_j|+\lambda_j)-C$.  The logarithm of the entire
ideal is the maximum of these expressions, since all coefficients have
absolute value at most one.  Positive multiplication and translation
commute with this nonempty finite maximum, giving $mND-C$.
The actual coefficient-interchange isomorphism and the intrinsic product
ideal comparison identify this whole ideal logarithm with the model
logarithm of the actual product ideal sheaf.
\end{proof}

\begin{lemma}[Making a finite model vertical]
\label{lem:pl-vertical-truncation}
\uses{lem:pl-fs-product-ideal}
\lean{EnvelopeProperty.ProjSubvariety.IsFiniteFS.exists_vertical_productIdeal_model}
\leanok
Let $\varphi$ be a rational finite Fubini--Study function on $X^{\rm an}$.
The ideal in the preceding lemma can be chosen to contain a positive power
of $t$, while preserving the identity
$\log|\mathcal J|(\sigma(v))=r\varphi(v)-C$ on $X^{\rm an}$, for an integer
$r>0$ and an integer $C$.
\end{lemma}
\begin{proof}
The continuous real function $r\varphi-C$ is bounded below on the compact
space $X^{\rm an}$.  Choose an integer $A>0$ such that
$-A\le r\varphi(v)-C$ everywhere.  Adjoin $t^A$ to the preceding finite
homogeneous generating set.  The logarithm of the enlarged ideal is
$\max\{-A,r\varphi-C\}=r\varphi-C$.  The added generator has $x$-degree
zero, so the enlarged ideal is still homogeneous and finitely generated.
Its inclusion of $t^A$ is an explicit algebraic central-support condition.
\end{proof}

\begin{proposition}[Vertical ideal data for piecewise linear functions]
\label{prop:pl-vertical-ideal-models}
\uses{lem:pl-vertical-truncation,lem:wu-locally-convex-pl}
\lean{EnvelopeProperty.ProjSubvariety.IsPL.exists_vertical_productIdeal_models,EnvelopeProperty.ProjSubvariety.IsLocallyConvex.exists_vertical_productIdeal_models}
\leanok
Every rational piecewise linear real function on $X^{\rm an}$ has the form
\[
 r\varphi(v)=\log|\mathcal I|(\sigma(v))
             -\log|\mathcal J|(\sigma(v))+c
\]
for an integer $r>0$, an integer $c$, and finitely generated homogeneous
product ideals $I,J$, each containing an explicit positive power of $t$.
The logarithms in this formula are finite on $X^{\rm an}$.  In particular
the assertion applies to a locally convex function in the sense of
Definition~\ref{def:wu-proj-functions}.
\end{proposition}
\begin{proof}
Write $\varphi=\alpha-\beta$ with rational finite Fubini--Study functions.
Apply the preceding lemma to obtain
$\log|\mathcal I_0|=r_1\alpha-C_1$ and
$\log|\mathcal J_0|=r_2\beta-C_2$.
Take $I=I_0^{r_2}$ and $J=J_0^{r_1}$.  Multiplicativity of intrinsic model
ideal norms gives
$\log|\mathcal I|=r_1r_2\alpha-r_2C_1$ and
$\log|\mathcal J|=r_1r_2\beta-r_1C_2$.
Subtract and take $r=r_1r_2$, $c=r_2C_1-r_1C_2$.
Both powers remain finitely generated, and contain respectively the
$r_2$-th and $r_1$-th powers of the original parameter generators.
The two displayed real-valued identities prove the finiteness required
for subtraction.  Local convexity supplies piecewise linearity by the
proved PL sheaf argument.
\end{proof}

\begin{remark}[Relation to local positivity]
\label{rmk:pl-ideals-remaining-geometry}
\uses{prop:pl-vertical-ideal-models,thm:wu-bj-testconfig,thm:wu-bj-nef-psh,thm:wu-cld-nef}
The construction produces product ideal sheaves and their intrinsic
Gauss model functions.  Restricting these ideals to the
$X\times\mathbf A^1$ model, the normal projective modification principalizing
them, its Cartier line bundle and the equality with that line
bundle's model function are the further geometric steps of
Theorem~\ref{thm:wu-bj-testconfig}.  Nonnegative degrees on central-fibre curves
and relative-nef Fubini--Study approximation are steps of
Theorems~\ref{thm:wu-cld-nef} and~\ref{thm:wu-bj-nef-psh}.  The proposition itself
does not assert local positivity.
\end{remark}
